% ═══════════════════════════════════════════════════════════════════════
\section{Un catalogue d'applications affines}
% ═══════════════════════════════════════════════════════════════════════

\begin{frame}{Les briques du chapitre}
  \small On passe en revue les applications affines élémentaires ---
  toutes des transformations, sauf la dernière (la projection). Ce sont
  les briques : toute application linéaire inversible du plan s'obtient
  en les composant (théorème de décomposition, plus loin).

  \vspace{6pt}
  Convention de notation, valable pour tout le chapitre : rotations en
  $r$ minuscule, réflexions en $R$ majuscule. $r_\theta$ / $R_\theta$ :
  centrées à l'origine. $r_{M,\theta}$ / $R_d$ : centre ou axe précisé
  explicitement --- \textbf{un indice à deux éléments}.
\end{frame}

\subsection{Translations}
\begin{frame}{Définition 1.5 --- Translation}
  \begin{Enonce}{Définition 1.5 --- Translation}
    La \defi{translation} de vecteur $\vec t = \col{e}{f}$ est
    \[
      t_{\vec t}\col{x}{y} = \col{x}{y} + \col{e}{f}
      = \mat{1}{0}{0}{1}\col{x}{y} + \col{e}{f}.
    \]
  \end{Enonce}
  \vspace{6pt}
  \small Partie linéaire = identité : aires et orientation préservées
  --- première isométrie. Composer deux translations donne une
  translation, et \textbf{l'ordre n'y change rien} :
  $t_{\vec u}\circ t_{\vec v} = t_{\vec u+\vec v} = t_{\vec v}\circ t_{\vec u}$.
\end{frame}

\subsection{Homothéties}
\begin{frame}{Définition 1.6 --- Homothétie de centre l'origine}
  \begin{Enonce}{Définition 1.6}
    L'\defi{homothétie} de centre $\vec 0$, rapport $c\in\Rs$ :
    $h_c\col{x}{y} = \mat{c}{0}{0}{c}\col{x}{y}$. $0<c<1$ : contraction
    uniforme ; $c>1$ : dilatation uniforme ; $c<0$ : retourne la figure
    ($h_{-1}$ = symétrie centrale).
  \end{Enonce}
  \small $\dete(cI) = c^2 > 0$ : une homothétie ne renverse
  \textbf{jamais} l'orientation, même si $c<0$ --- excellent sujet de
  discussion en classe.
\end{frame}

\begin{frame}{Proposition 1.5 --- Homothétie de centre quelconque}
  \begin{Enonce}{Proposition 1.5}
    L'homothétie de centre $M=(h,k)$, rapport $c$ :
    $h_{M,c} = t_{\vec m}\circ h_c \circ t_{-\vec m}$,
    \[
      h_{M,c}\col{x}{y} = \col{cx+h(1-c)}{cy+k(1-c)}.
    \]
  \end{Enonce}
  \begin{AuTableau}
    \small Technique de la \textbf{conjugaison} (sert partout) : on
    ramène le centre à l'origine, on fait la transformation, on
    revient. $P\mapsto P-\vec m\mapsto c(P-\vec m)\mapsto c(P-\vec
    m)+\vec m$.
  \end{AuTableau}
\end{frame}

\begin{frame}{Pour la classe 1.3 --- Lire une formule du secondaire}
  \begin{Classe}{Lire une formule du secondaire}
    \small $(x,y)\mapsto(cx+h(1-c), cy+k(1-c))$ --- exactement la
    formule de 2\up{e} secondaire. Vue comme
    $t_{\vec m}\circ h_c\circ t_{-\vec m}$ : « je vais au centre,
    j'agrandis, je reviens » --- rien à retenir.
  \end{Classe}
\end{frame}

\subsection{Rotations}
\begin{frame}{Définition 1.7 --- Rotation de centre l'origine}
  \begin{Enonce}{Définition 1.7}
    La \defi{rotation} d'angle $\theta$, sens antihoraire, autour de
    l'origine :
    $r_\theta = \mat{\cos\theta}{-\sin\theta}{\sin\theta}{\cos\theta}$.
  \end{Enonce}
  \FigBox[0.42\textheight]{%
    \begin{tikzpicture}[scale=1.4]
      \repere{-1.35}{1.6}{-1.35}{1.6}
      \draw[gray!45, thin] (0,0) circle (1);
      \draw[fleche] (0,0) -- (1,0) node[etiq, below right] {$\vec{e}_1$};
      \draw[fleche] (0,0) -- (0,1) node[etiq, below left] {$\vec{e}_2$};
      \draw[fleche, firebrick3] (0,0) -- (37:1)
        node[etiq, right] {$\col{\cos\theta}{\sin\theta}$};
      \draw[fleche, firebrick3] (0,0) -- (127:1)
        node[etiq, above left] {$\col{-\sin\theta}{\cos\theta}$};
      \draw[firebrick3, thin] (0.42,0) arc (0:37:0.42);
      \node[etiq, firebrick3] at (20:0.62) {$\theta$};
      \draw[firebrick3, thin] (0,0.42) arc (90:127:0.42);
    \end{tikzpicture}}
  \begin{center}\small Figure 1.5. $\dete(r_\theta)=1$ : aires et
    orientation préservées.\end{center}
\end{frame}

\begin{frame}{Remarque 1.7 --- Un ou deux indices ?}
  \begin{Rem}{Remarque 1.7}
    Un indice \textbf{seul}, comme $r_\theta$ : centrée à
    \textbf{l'origine}. Dès que le centre n'est pas l'origine, on le
    place en \textbf{premier} indice : $r_{M,\theta}$ --- l'indice
    porte alors sur \emph{deux} objets, jamais sur l'angle seul. Même
    lecture pour $h_c$ / $h_{M,c}$ et $R_\theta$ / $R_d$.
  \end{Rem}
  \[
    r_{M,\theta}\col{x}{y}
      = r_\theta\left(\col{x}{y} - \vec{m}\right) + \vec{m}
      = r_\theta\col{x}{y} + (I - r_\theta)\,\vec{m}
      \qquad\text{(conjugaison, encore.)}
  \]
\end{frame}

\begin{frame}{Proposition 1.6 --- Composer deux rotations de même centre}
  \begin{Enonce}{Proposition 1.6}
    Pour $M$ quelconque :
    $r_{M,\alpha}\,r_{M,\beta} = r_{M,\alpha+\beta} = r_{M,\beta}\,r_{M,\alpha}$,
    et $r_{M,\theta}^{-1} = r_{M,-\theta}$.
  \end{Enonce}
  \begin{AuTableau}
    \small Idée : cas $M=\vec 0$ par produit matriciel (les formules
    d'addition d'angles apparaissent) ; cas général par conjugaison,
    ramené au premier cas.
  \end{AuTableau}
  \begin{Classe}{Un cadeau pour la classe}
    \small La démonstration (cas $M=\vec 0$) est une preuve de
    $\cos(\alpha+\beta)=\cos\alpha\cos\beta-\sin\alpha\sin\beta$ et
    $\sin(\alpha+\beta)=\sin\alpha\cos\beta+\cos\alpha\sin\beta$ --- la
    plus courte qui soit, et elle explique \emph{pourquoi} : reflet
    d'un produit de matrices. À garder pour la 5\up{e} secondaire.
  \end{Classe}
\end{frame}

\subsection{Réflexions}
\begin{frame}{Définition 1.8 --- Réflexions}
  \begin{Enonce}{Définition 1.8 --- Réflexions par rapport aux axes}
    $H = \mat{1}{0}{0}{-1}$ (axe des $x$), $V=\mat{-1}{0}{0}{1}$ (axe
    des $y$).
  \end{Enonce}
  Plus généralement, réflexion d'axe passant par l'origine, angle
  $\theta$ avec l'axe des $x$ :
  \[
    R_\theta = \mat{\cos 2\theta}{\sin 2\theta}{\sin 2\theta}{-\cos 2\theta}.
  \]
  \small $\dete(R_\theta) = -1$ : une réflexion renverse toujours
  l'orientation (intuition du miroir). $\theta=0$ : $H$. $\theta=\pi/4$ :
  réflexion d'axe $y=x$, matrice $\mat{0}{1}{1}{0}$ (échange les
  coordonnées).
\end{frame}

\begin{frame}{Remarque 1.8 --- L'angle double}
  \begin{Rem}{Remarque 1.8}
    Le $2\theta$ n'est \textbf{pas} une coquille : faire tourner le
    \emph{miroir} de $\theta$ fait tourner l'\emph{image} de
    $2\theta$. Même phénomène : la composée de deux réflexions d'axes
    sécants est une rotation d'angle \textbf{double} de l'angle des
    axes (théorème~1.13, plus loin) --- l'un des plus beaux résultats
    à faire découvrir au secondaire.
  \end{Rem}
\end{frame}

\subsection{Cisaillements et contractions non uniformes}
\begin{frame}{Définition 1.9 --- Cisaillement}
  \begin{Enonce}{Définition 1.9}
    \defi{Cisaillement horizontal} de rapport $m$ : $\mat{1}{m}{0}{1}$.
    \defi{Cisaillement vertical} de rapport $m$ : $\mat{1}{0}{m}{1}$.
  \end{Enonce}
  \small Le cisaillement horizontal laisse fixe l'axe des $x$ et fait
  glisser chaque droite horizontale proportionnellement à sa hauteur ---
  un paquet de cartes qui penche. $\dete = 1$ : \textbf{un cisaillement
  conserve les aires}, en déformant complètement angles et longueurs.
  \begin{Classe}{Le cisaillement et l'aire du parallélogramme}
    \small « Aire = base $\times$ hauteur » : un cisaillement bien
    choisi transforme le parallélogramme en un rectangle de même base
    et hauteur, en conservant l'aire. C'est exactement ce que l'on fait
    au secondaire en découpant un triangle et en le déplaçant --- sans
    savoir qu'on vient de faire un cisaillement.
  \end{Classe}
\end{frame}

\begin{frame}{Définition 1.10 --- Contraction non uniforme}
  \begin{Enonce}{Définition 1.10}
    Application linéaire $\mat{a}{0}{0}{b}$, $0<a<1$, $0<b<1$.
  \end{Enonce}
  \small Si $a\neq b$ : contracte davantage dans une direction ---
  un cercle devient une ellipse. Avec $a>1$, $0<b<1$ : dilatation
  horizontale et contraction verticale. Essentielles dans la
  construction des fractales \emph{auto-affines} (fougère de Barnsley).
\end{frame}

\begin{frame}{Figure 1.6 --- Le catalogue en images}
  \FigBox[0.66\textheight]{%
    \begin{tikzpicture}[scale=0.72]
      \begin{scope}[shift={(0,0)}]
        \draw[quadrillage] (-1.25,-1.25) grid (1.25,1.25);
        \draw[axe] (-1.25,0) -- (1.4,0);
        \draw[axe] (0,-1.25) -- (0,1.4);
        \draw[steelblue4!55, thin] \motifF;
        \begin{scope}[cm={1,0,0,1,(-0.9,-0.6)}]
          \fill[figim] \motifF;
        \end{scope}
        \node[etiq, black, align=center] at (0,-1.72)
          {translation\\$\vec{t}=(-0{,}9\,;-0{,}6)$};
      \end{scope}
      \begin{scope}[shift={(3.5,0)}]
        \draw[quadrillage] (-1.25,-1.25) grid (1.25,1.25);
        \draw[axe] (-1.25,0) -- (1.4,0);
        \draw[axe] (0,-1.25) -- (0,1.4);
        \draw[steelblue4!55, thin] \motifF;
        \begin{scope}[cm={0,1,-1,0,(0,0)}]
          \fill[figim] \motifF;
        \end{scope}
        \node[etiq, black, align=center] at (0,-1.72)
          {rotation\\$\theta = 90°$};
      \end{scope}
      \begin{scope}[shift={(7,0)}]
        \draw[quadrillage] (-1.25,-1.25) grid (1.25,1.25);
        \draw[axe] (-1.25,0) -- (1.4,0);
        \draw[axe] (0,-1.25) -- (0,1.4);
        \draw[steelblue4!55, thin] \motifF;
        \begin{scope}[cm={1,0,0,-1,(0,0)}]
          \fill[figim] \motifF;
        \end{scope}
        \node[etiq, black, align=center] at (0,-1.72)
          {réflexion\\d'axe $Ox$};
      \end{scope}
      \begin{scope}[shift={(0,-4.3)}]
        \draw[quadrillage] (-1.25,-1.25) grid (1.25,1.25);
        \draw[axe] (-1.25,0) -- (1.4,0);
        \draw[axe] (0,-1.25) -- (0,1.4);
        \draw[steelblue4!55, thin] \motifF;
        \begin{scope}[cm={0.55,0,0,0.55,(0,0)}]
          \fill[figim] \motifF;
        \end{scope}
        \node[etiq, black, align=center] at (0,-1.72)
          {homothétie\\$c = 0{,}55$};
      \end{scope}
      \begin{scope}[shift={(3.5,-4.3)}]
        \draw[quadrillage] (-1.25,-1.25) grid (1.25,1.25);
        \draw[axe] (-1.25,0) -- (1.4,0);
        \draw[axe] (0,-1.25) -- (0,1.4);
        \draw[steelblue4!55, thin] \motifF;
        \begin{scope}[cm={1,0.8,0,1,(0,0)}]
          \fill[figim] \motifF;
        \end{scope}
        \node[etiq, black, align=center] at (0,-1.72)
          {cisaillement\\vertical, $m=0{,}8$};
      \end{scope}
      \begin{scope}[shift={(7,-4.3)}]
        \draw[quadrillage] (-1.25,-1.25) grid (1.25,1.25);
        \draw[axe] (-1.25,0) -- (1.4,0);
        \draw[axe] (0,-1.25) -- (0,1.4);
        \draw[steelblue4!55, thin] \motifF;
        \begin{scope}[cm={0.5,0,0,0.85,(0,0)}]
          \fill[figim] \motifF;
        \end{scope}
        \node[etiq, black, align=center] at (0,-1.72)
          {contraction\\non uniforme};
      \end{scope}
    \end{tikzpicture}}
  \begin{center}\small Le « F » : seule la réflexion renverse
  l'orientation.\end{center}
\end{frame}

\begin{frame}{Table 1.1 --- Récapitulatif}
  \small
  \begin{center}
  \begin{tabular}{llcc}
    \toprule
    Transformation & Matrice $A$ & $\dete(A)$ & Isométrie ? \\
    \midrule
    Translation & $I$ & $1$ & oui \\[4pt]
    Rotation d'angle $\theta$
      & $\mat{\cos\theta}{-\sin\theta}{\sin\theta}{\cos\theta}$
      & $1$ & oui \\[10pt]
    Réflexion d'axe $\theta$
      & $\mat{\cos 2\theta}{\sin 2\theta}{\sin 2\theta}{-\cos 2\theta}$
      & $-1$ & oui \\[10pt]
    Homothétie de rapport $c$
      & $\mat{c}{0}{0}{c}$ & $c^2$ & si $\lvert c\rvert = 1$ \\[10pt]
    Cisaillement de rapport $m$
      & $\mat{1}{m}{0}{1}$ & $1$ & non (sauf $m=0$) \\[10pt]
    Contraction non uniforme
      & $\mat{a}{0}{0}{b}$ & $ab$ & non (sauf $\lvert a\rvert=\lvert b\rvert=1$) \\[10pt]
    Projection sur $y=x$
      & $\tfrac12\mat{1}{1}{1}{1}$ & $0$ & non (non inversible) \\
    \bottomrule
  \end{tabular}
  \end{center}
\end{frame}

\subsection{Le cas non inversible}
\begin{frame}{Le cas non inversible}
  \small On considère parfois des applications affines qui ne sont pas
  des transformations : $\dete(A) = 0$. Tout le plan est alors écrasé
  sur une droite ($A\neq 0$) ou un point ($A=0$).

  \begin{Enonce}{Exemple 1.7 --- La projection orthogonale sur $y=x$}
    \small $p\col{x}{y} = \tfrac{x+y}{2}\col{1}{1}
      = \tfrac12\mat{1}{1}{1}{1}\col{x}{y}$. $\dete = 0$ : toutes les
      aires sont écrasées à zéro. Notez $p\circ p = p$.
  \end{Enonce}
  \FigBox[0.32\textheight]{%
    \begin{tikzpicture}[scale=1.15]
      \repere{-1.4}{2.4}{-1.4}{2.4}
      \draw[droite] (-1.3,-1.3) -- (2.3,2.3) node[etiq, right, black] {$y=x$};
      \fill[fig] \motifF;
      \draw[firebrick3, ultra thick] (0,0) -- (0.875,0.875);
      \foreach \x/\y in {0/1, 0.75/1, 0.3/0, 0.6/0.55}
        {\draw[gray!70, thin, dotted] (\x,\y) -- ({(\x+\y)/2},{(\x+\y)/2});
         \node[pt, inner sep=0.7pt] at (\x,\y) {};
         \node[ptim, inner sep=0.7pt] at ({(\x+\y)/2},{(\x+\y)/2}) {};}
    \end{tikzpicture}}
  \begin{center}\small Figure 1.7 --- Le motif tout entier écrasé sur un
  segment. Pas une transformation du plan.\end{center}
\end{frame}

\subsection{Le théorème de décomposition}
\begin{frame}{Théorème 1.7 --- Décomposition}
  \begin{Enonce}{Théorème 1.7 --- Décomposition des applications linéaires}
    Une application linéaire de $\R^2$ est inversible si et seulement
    si elle est une composition d'homothéties, de rotations, de
    réflexions, de cisaillements et de contractions/dilatations non
    uniformes.
  \end{Enonce}
  \small C'est la version géométrique d'un fait d'algèbre linéaire :
  toute matrice inversible est un produit de matrices élémentaires.
  Cisaillements $\leftrightarrow$ addition d'un multiple d'une ligne à
  une autre ; contractions non uniformes $\leftrightarrow$
  multiplication d'une ligne par un scalaire ; réflexion d'axe $y=x$
  $\leftrightarrow$ échange de deux lignes. \textbf{Gauss--Jordan est
  une décomposition géométrique.}
  \begin{Rem}{Remarque 1.9 --- Deux décompositions modernes}
    \footnotesize \emph{Polaire} : $A=QS$ (rotation/réflexion, puis
    déformation symétrique positive). \emph{SVD} : $A=U\Sigma V^\top$
    --- « une rotation, un étirement selon les axes, une rotation ».
    Explique pourquoi un cercle devient toujours une ellipse.
  \end{Rem}
\end{frame}
